\subsection{The Instability of the Subject: The Inherent Challenge of Evaluating LLMs}

Hallucination represents a significant challenge for contemporary generative large language models \cite{huang2025survey}. With regard to general natural language tasks, there has been some exploration of LLM hallucinations to a certain extent \cite{tonmoy2024comprehensive}\cite{ye2023cognitive}\cite{zhang2025siren}, and these are typically categorised into three types: 'Input-Conflicting Hallucination', 'Fact-Conflicting Hallucination' and 'Context-Conflicting Hallucination' \cite{zhang2025siren}.In specific domains, such hallucination pose direct and severe risks, representing one of the most significant challenges facing LLMs.

The issue of factuality is further complicated by the challenge of inconsistent and non-deterministic outputs\cite{atil2024non}\cite{stureborg2024large}. It is widely acknowledged that LLMs demonstrate a high degree of sensitivity to input phrasing\cite{an2024make}, minor alterations in system prompts\cite{he2024does}, and even decoding parameters\cite{spector2023accelerating}. This sensitivity leads to significantly varied responses to questions with identical semantics\cite{atil2024non}. This absence of determinism suggests that a model may provide a correct answer at one moment, a different answer at another, or deliver contradictory reasoning processes across multiple attempts\cite{turpin2023language}. Such behaviour undermines the core expectation of robustness in software systems and renders performance evaluation highly dependent on often arbitrary testing configurations\cite{chen2024chatunitest}. It is imperative that an evaluation framework is implemented in order to verify the correctness of individual attempts, whilst also measuring the stability and reproducibility of a model's outputs across a range of prompts and conditions.

In order to address the knowledge gaps and 'hallucination' tendencies inherent in foundational large language models, Retrieval-Augmented Generation (RAG) has emerged as a mainstream architectural paradigm\cite{liu2024much}. By anchoring the generation process to externally retrieved documents, RAG aims to improve both factual accuracy and controllability\cite{gao2023retrieval}\cite{guu2020retrieval}. However, this introduces new layers of complexity and cascading vulnerabilities into the system\cite{agrawal2024mindful}. RAG pipeline performance is no longer solely dependent on a single model, but critically depends on the coordinated behaviour of its constituent components: the retriever, the reordering module , and the generator\cite{gao2023retrieval}\cite{yan2024corrective}.
In some survey, the retriever may fail to obtain the most relevant documents or retrieve outdated or contradictory information\cite{wang2025retrieval}. Even with perfect retrieval, the generator may disregard the provided context, synthesise information erroneously from multiple sources, or produce speculative content that is not supported by the evidence \cite{yu2024evaluation}. This complexity transforms the evaluation task from judging a single model's output to diagnosing a multi-component system where error sources are opaque and interdependent\cite{wang2024evaluating}. Consequently, an effective evaluation framework must possess sufficient granularity to determine whether errors originate from knowledge retrieval, information synthesis, or adherence to the basis\cite{es2024ragas}. The fragility of such augmented systems indicates that the 'object' of evaluation is often a fragile composite. This means that evaluation tools must be robust enough to address and diagnose this complexity without introducing significant additional noise\cite{ni2025towards}\cite{yu2024evaluation}.

\subsection{2.3 The Instability : Critical Flaws in Prevailing Evaluation Methodologies}

In order to address the challenges inherent in the evaluation of open-domain generation tasks, there has been a rapid rise in the utilisation of more powerful large language models (e.g. GPT-4) as "judges" to assess the outputs of other models\cite{li2025generation}\cite{zheng2023judging}\cite{szymanski2025limitations}. While this approach is appealing due to its flexibility and its use of an anthropomorphic judgement, it need to be noted that it shifts the core risk of evaluation from the assessed model to another evaluator model. This introduces unacceptable levels of variance and manipulability, thereby undermining the very foundations of objective benchmarking\cite{wataoka2024self}\cite{wei2024systematic}.

Research consistently demonstrates that such "LLM-as-a-Judge" outputs exhibit extreme sensitivity to the minutiae of prompt wording, the format of scoring criteria definitions, and even the tone of system instructions\cite{li2025generation}\cite{zheng2023judging}\cite{szymanski2025limitations}. This phenomenon gives rise to substantial variations in the scoring of identical responses when assessed across different prompts\cite{he2024does}. This sensitivity is not merely a minor error but rather a systemic bias, which renders the evaluation results difficult to reproduce and heavily dependent on the specific configuration of the evaluator model\cite{perlitz2024benchmark}\cite{banerjee2024vulnerability}. Of greater concern is the fact that this paradigm permits systematic manipulation\cite{zhou2023don}. The meticulous crafting of evaluation prompts by assessors can result in the steering of these models towards scoring biases for specific output types\cite{dietz2025llm}. This can lead to the artificial inflation of benchmark scores and the undermining of the fundamental purpose of fair comparison. The model under scrutiny does not serve as the primary metric in this instance; rather, the focus is directed towards the proficiency of 'prompt engineering'\cite{bruni2025benchmarking}\cite{liang2025much}. It can be argued that, rather than resolving reliability issues, this evaluation method—due to its inherent instability and manipulability—becomes a new, more insidious source of noise\cite{feuer2025judgment}. This approach is in direct opposition to the fundamental principles of objectivity, stability, and reproducibility, which are prerequisites for the establishment of a reliable benchmark.

In tasks requiring definitive answers, traditional aggregation metrics such as F1 score and exact match have long been regarded as the gold standard\cite{hu2024unveiling}\cite{wu2024detectrl}. However, when evaluating the complex generated content of large language models, these metrics reveal fundamental limitations: calculated through superficial string matching, they fail to capture semantic fidelity, logical reasoning, or determine whether a ‘correct’ answer stems from a sound reasoning process\cite{white2024livebench}\cite{qiu2024efficient}. Models may yield matching answers through flawed logic or coincidental data patterns, rendering current metrics entirely ineffective in such scenarios\cite{huang2023large}\cite{zeng2023evaluating}.
More critically, when applied to open-ended generation assessments, these traditional metrics are seldom computed directly\cite{cheng2025survey}. Instead, they rely on an unstable upstream judgement\cite{xu2024benchmark}. For instance, to calculate similarity between a response and a reference answer, the system may first invoke another LLM to determine semantic equivalence between the two. Within this chain, the seemingly objective numerical results of traditional metrics actually inherit all the instability, variance, and bias inherent in the upstream ‘LLM-as-a-Judge’ evaluator\cite{li2024llms}\cite{lee2025evaluating}. As demonstrated in the extant literature, it may be concluded that the aggregation of metrics incorporates unstable probabilistic judgements into superficially precise point estimates, thereby masking rather than eliminating the fundamental uncertainty inherent in the evaluation. Consequently, the final benchmark score is a noisy signal, with fluctuations primarily reflecting the randomness of the evaluation chain rather than genuine differences in model performance.

In summary, the two prevailing evaluation paradigms share a fundamental flaw: they both fail to resolve, and indeed exacerbate, the issue. The employment of unstable generative scorers, whether directly or indirectly through the implementation of traditional metrics, invariably gives rise to an unstable measurement. From a statistical perspective, this results in the confidence intervals for model performance estimates being excessively broad, thereby rendering subtle performance differences between models statistically insignificant. Consequently, the application of nuanced comparisons and rankings becomes redundant.

\subsection{The Specific Challenge of Evaluating Memory and Knowledge}
In recent years, the research community has come to acknowledge the significance of specialised evaluations. Numerous benchmarks have been proposed with the aim of assessing long-term contextual memory, factual accuracy, and knowledge attribution\cite{liu2023think}\cite{li2024enhancingllmfactualaccuracy}\cite{wang2024should}. These efforts have yielded substantial contributions towards constructing high-quality, challenging test datasets, designing complex multi-hop reasoning problems, tasks requiring precise recall of lengthy documents, and adversarial samples for testing factual consistency\cite{gekhman2023trueteacher}\cite{tam2023evaluating}.
However, concomitant with these advances, a fundamental methodological stagnation has emerged. Notwithstanding the highly specialised nature of the test content, the prevailing scoring methods continue to be anchored within the unstable paradigm that was the subject of criticism in Section 2.3.
Consequently, when evaluating the capacity of a model to retain or forget information, we depend on automated evaluators that may themselves demonstrate inconsistencies or biases in their recall\cite{abeysinghe2024challenges}\cite{chen2025multi}. Notwithstanding the pressing need for memory- and knowledge-oriented evaluation, a framework designed from the ground up to measure this capability stably and reliably remains methodologically absent. This scenario represents more than a mere application of extant benchmarks; it is the domain in which the inherent flaws of these benchmarks are most starkly exposed. This issue is of particular significance as it represents a compelling call for an entirely new, stable evaluation paradigm.

\subsection{Identifying the Gap for a Probabilistically Stable Benchmark}
The present chapter's discussion reveals a profound core contradiction within the field of large language model evaluation: an attempt is being made to measure an inherently unstable object using a series of tools that are themselves unstable. On the one hand, LLMs and their augmented systems possess inherent deficiencies in factuality, consistency, and robustness (Section 2.2); on the other hand, the two predominant approaches dominating the current evaluation ecosystem – generative LLM scoring and indirect assessment based on traditional metrics – are both widely criticised for exhibiting excessive "evaluator variance". Rather than providing definitive measurements, they yield highly variable probability distributions that are susceptible to manipulation (see Section 2.3). This contradiction has precipitated a particularly severe credibility crisis in the assessment of critical capabilities such as memory and knowledge (see Section 2.4).
A critical synthesis of extant literature indicates that the most pressing challenge facing the field is no longer merely the need for more and harder test items, but the fragility of the assessment scientific infrastructure. The present research deficit is not a paucity of data, but rather a dearth of methodology. Specifically, a benchmarking framework that prioritises 'probabilistic stability' and 'measurement reliability' as its foremost design objectives is lacking. In order to address this gap, it is necessary to establish a framework that possesses the following core characteristics:
The property of low variance is complemented by the property of high reproducibility. The scoring mechanism should yield statistically stable results under identical conditions, thereby minimising random noise introduced by the evaluation process itself.
The capacity to resist manipulation is a pivotal factor in the study of human behaviour. The logic underpinning its assessment must demonstrate robustness against superficial optimisation strategies, such as prompt engineering, thereby ensuring that scores accurately reflect a model's intrinsic capabilities rather than exploiting evaluation vulnerabilities.
The requirement for transparency and interpretability is paramount. The process of scoring should be designed in such a way as to minimise reliance on complex black-box generative adjudicators. Instead, the focus should be on the use of more deterministic and analysable computational methods. This approach is intended to ensure the objectivity of the evaluation process.
Consequently, a genuinely urgent breakthrough is required, and this lies in shifting from 'pursuing smarter evaluators' to 'designing more stable evaluators'. The fundamental methodological gap that this research seeks to address constitutes the motivation for this study. The MemIndex benchmark proposed in this paper is an exemplar of this approach in practice. The present study posits that the test set under scrutiny is not merely another memory test set; its core contribution lies in introducing an entirely new evaluation paradigm. This paradigm is grounded in probabilistically stable binary classification, with the aim of providing a low-variance, reproducible, and manipulation-resistant metric for measuring memory capability. By systematically demonstrating the instability of traditional approaches and validating the robustness of the new paradigm, MemIndex seeks to establish a more solid and credible scientific foundation for evaluating LLM memory capabilities.
